% Source pédagogique : cours/cm1.qmd. Choc d'offre, élasticité, répartition.
\begin{board}{1}{Ormuz : le choc passe surtout dans le prix}
 {La hausse pèse davantage sur les ménages dépendants de la voiture.}
\stage{5.4}{1}{Moins de pétrole}
\stage{16}{2}{Des trajets\\peu ajustables}
\stage{26.6}{3}{Le prix réagit fortement}
% Détroit : deux rives et un passage resserré.
\fill[blue!7] (1.5,7) rectangle (9.3,12.6);
\draw[drawline,fill=ink!7] (1.5,12.6)--(4.3,12.6)--(4.9,10.5)--(4.2,9.8)--(4.2,7)--(1.5,7);
\draw[drawline,fill=ink!7] (9.3,12.6)--(6.6,12.6)--(6,10.5)--(6.7,9.8)--(6.7,7)--(9.3,7);
\draw[amber,line width=5pt] (4.9,10.5)--(5.75,10.5);
\draw[drawline,fill=paper] (5,8.9)--(6,8.9)--(5.8,8.5)--(5.2,8.5)--cycle;
\draw[drawline,fill=paper] (5.3,8.9) rectangle (5.7,9.3);
\draw[flow] (5.45,8)--(5.45,6.8);
\captext{5.4}{Choc d'offre négatif :\\moins de barils disponibles.}
\house{14.1}{9.1}
\car{16.5}{6.7}
\draw[drawline] (18.2,10)--(19.8,10)--(19.8,12.1)--(18.2,12.1)--cycle;
\draw[blue,line width=2pt] (18.5,10.5)--(19.5,10.5) (18.5,11.6)--(19.5,11.6);
\draw[amber,line width=2.5pt] (18,9.8)--(20,12.3);
\captext{16}{À court terme :\\demande peu élastique.}
\pump{24.5}{8.1}
\say{28.3}{11}{$P\ \uparrow$}
\barrel{26.6}{6.8}
\say{29.1}{7.4}{$Q\ \downarrow$}
\captext{26.6}{Le prix monte beaucoup ;\\la quantité baisse peu.}
\end{board}
